\section*{Tóm tắt đề án}

Báo cáo đề án tập trung giải quyết bài toán quản trị chất lượng thực tiễn tại Trường Cao đẳng nghề Kỹ thuật -- Công nghiệp A, một cơ sở giáo dục nghề nghiệp có quy mô 3.200 học sinh, sinh viên và đội ngũ 180 nhà giáo, trong bối cảnh nhà trường vừa hoàn thành đợt tự đánh giá cấp cơ sở năm 2026 nhưng chỉ đạt 79 trên tổng số 100 điểm, không đủ điều kiện công nhận đạt tiêu chuẩn kiểm định chất lượng theo quy định hiện hành. Đề án tiếp cận theo hướng chẩn đoán nguyên nhân gốc rễ, thiết lập bộ công cụ tự đánh giá chuẩn mực cho chương trình đào tạo mũi nhọn trình độ cao đẳng nghề Điện công nghiệp và tái thiết kế hệ thống đảm bảo chất lượng bên trong gắn với kế hoạch hành động cải tiến chất lượng chu kỳ 01 năm.

Về phương diện chẩn đoán pháp lý, đề án tiến hành đối chiếu định lượng kết quả tự đánh giá cấp cơ sở của nhà trường với quy định tại Khoản 1 Điều 5 Thông tư số 14/2024/TT-BLĐTBXH, khẳng định nhà trường vi phạm đồng thời cả hai điều kiện bắt buộc: vi phạm Điều kiện a do tổng điểm không đạt ngưỡng sàn 80 điểm và vi phạm Điều kiện c do cả hai tiêu chí trọng tâm là Tiêu chí 2 về Hoạt động đào tạo và Tiêu chí 5 về Cơ sở vật chất, thiết bị đào tạo đều chỉ đạt 75.0\% điểm chuẩn, dưới ngưỡng sàn 80\%. Vận dụng các quan điểm Giá trị gia tăng, Phù hợp mục tiêu, Thỏa mãn khách hàng và quy tắc 4M, đề án làm rõ nguyên nhân bắt nguồn từ sự thụ động của bộ máy quản trị chất lượng (Ban Đảm bảo chất lượng mỗi năm chỉ họp một lần), sự quá tải và đơn độc của Phòng Đào tạo trong lưu trữ minh chứng, tâm lý e ngại cung cấp số liệu của các đơn vị và sự thiếu đồng bộ của trang thiết bị xưởng thực hành.

Về phương diện thực hành tự đánh giá chương trình đào tạo, đề án vận dụng Chương III và Phụ lục 03 của Thông tư số 14/2024/TT-BLĐTBXH để thiết lập bộ tiêu chí gồm 07 tiêu chí, 40 tiêu chuẩn (thang điểm 100) và khung danh mục 40 minh chứng chuẩn hóa cho nghề Điện công nghiệp; tiến hành thẩm định thử 05 bộ hồ sơ minh chứng điển hình; biên soạn trọn vẹn Báo cáo tự đánh giá cho Tiêu chuẩn 1.3 về tích hợp năng lực số; và xác lập bảng điểm dự kiến đạt 82.5 điểm, đáp ứng đầy đủ cả ba điều kiện công nhận đạt chuẩn tại Điều 11 của thông tư kiểm định.

Về phương diện giải pháp hệ thống, đề án vận dụng mô hình 05 bước của David Lim (2001) để tái cấu trúc hệ thống đảm bảo chất lượng bên trong, xóa bỏ cơ chế kiêm nhiệm hình thức, thiết lập mạng lưới đảm bảo chất lượng tại từng khoa và ban hành quy chế chia sẻ dữ liệu nội bộ. Đồng thời, đề án xây dựng Kế hoạch hành động cải tiến chất lượng chu kỳ 01 năm theo chu trình Plan -- Do -- Check -- Act (PDCA), thiết lập ma trận phân công trách nhiệm RACI và phân bổ gói ngân sách 3.5 tỷ đồng nhằm đưa Tiêu chí 2 và Tiêu chí 5 vượt ngưỡng 80\% điểm chuẩn, bảo đảm nhà trường hoàn thành tự đánh giá lại cấp cơ sở và tự tin đón nhận đoàn đánh giá ngoài đối với chương trình đào tạo nghề Điện công nghiệp.

\vspace{0.5cm}
\noindent\textbf{Từ khóa:} Quản lý chất lượng giáo dục nghề nghiệp; Thông tư số 14/2024/TT-BLĐTBXH; Đảm bảo chất lượng bên trong (IQA); Chu trình PDCA; Nghề Điện công nghiệp; Đánh giá ngoài; Quy tắc 4M.
